% generated by tools/gendoc.py from the doc comments in src/; do not edit
\fn{(fft signal)}{=\textgreater{} (list re im), each of length next-pow2(length); the signal is zero-padded}
\fn{(ifft (list re im)}{) =\textgreater{} real signal of the same length (a power of two)}
\fn{(osc sr freqs table)}{=\textgreater{} one sample per element of freqs, reading the wavetable with linear}
\fn{(iir x b a)}{=\textgreater{} y with a0 y[n] = sum b[k] x[n-k] - sum a[k\textgreater{}0] y[n-k]; a[0] need not be 1}
\fn{(delay x d)}{=\textgreater{} x delayed by d samples (fractional, linear interpolation), same length}
\fn{(resample x factor)}{=\textgreater{} x with its length multiplied by factor, by windowed-sinc interpolation}
\fn{(autocorr x)}{=\textgreater{} biased autocorrelation (sum / n) for lags 0 .. n/2 - 1; it decays with lag, which}
\fn{(interleave (list a b ...)}{) =\textgreater{} one vector a0 b0 a1 b1 ...}
\fn{(deinterleave v channels)}{=\textgreater{} (list a b ...): the channels of an interleaved vector}
\fn{(local-maxima v)}{=\textgreater{} indices k with v[k] \textgreater{} v[k-1] and v[k] \textgreater{} v[k+1] (interior points only)}
\fn{(gather v indices)}{=\textgreater{} the elements of v at the given (integer) indices, as a vector}
\fn{(add-at! dst pos src)}{=\textgreater{} dst, with src added in place starting at pos; dst must be long enough.}
\fn{(comb x d g)}{=\textgreater{} feedback comb y[n] = x[n] + g y[n-d]; (allpass x d g) =\textgreater{} Schroeder allpass}
\fn{(pvoc x opts)}{after sparkle: fixed synthesis hop I = window/overlap, analysis hop D = I/stretch}
\fn{(sine sr freq dur)}{a sine of freq Hz lasting dur seconds}
\fn{(noise n)}{white noise in [-1, 1]}
\fn{(impulse n)}{a 1 followed by n-1 zeros}
\fn{(gen n amps)}{one period of a wavetable with the given harmonic amplitudes, plus a guard point (n+1 samples, last = first), normalised to a peak of 1; for osc}
\fn{(oscbank sr amps freqs table)}{sum of oscillators; amps and freqs are lists of per-sample envelopes}
\fn{(mix layers)}{overlay (list (list position signal) ...) into one signal}
\fn{(add-at dst pos src)}{dst with src added starting at pos; dst grows if needed}
\fn{(fade-in x n)  (fade-out x n)}{linear fades over n samples}
\fn{(rms x)}{root mean square}
\fn{(normalize-peak x)}{scale so the peak is 1; (normalize-rms x target) so the rms is target}
\fn{(normalize-rms x target)}{scale so the rms is target}
\fn{(db x)  (undb d)}{amplitude to decibels and back}
\fn{(next-pow2 n)}{the smallest power of two \textgreater{}= n (fft sizes)}
\fn{(window n a0 a1 a2)}{generalised cosine window, periodic (so hann overlap-adds exactly at hop n/4); (hann n) (hamming n) (blackman n) are the usual ones}
\fn{(hann n)  (hamming n)  (blackman n)}{the usual windows, periodic}
\fn{(hamming n)}{the Hamming window}
\fn{(blackman n)}{the Blackman window}
\fn{(car-\textgreater{}pol spec)}{(list re im) -\textgreater{} (list mag phase); (pol-\textgreater{}car) the other way}
\fn{(pol-\textgreater{}car spec)}{(list mag phase) -\textgreater{} (list re im)}
\fn{(magnitudes spec)  (phases spec)}{from a complex spectrum}
\fn{(spectrum-freqs n sr)}{the frequency of each of the first n/2 bins}
\fn{(magnitude-spectrum x)}{the positive-frequency magnitudes of x, normalised so a full-scale sine reads 1}
\fn{(complex-mul a b)}{product of two (list re im) spectra}
\fn{(conv x h)}{linear convolution through the FFT, length (+ (length x) (length h) -1)}
\fn{(cepstrum mags)}{real cepstrum of a full (symmetric) magnitude spectrum}
\fn{(spectral-envelope mags order)}{smooth envelope of a full magnitude spectrum: the "true envelope", a cepstral smoothing (first `order` coefficients) iterated so that it rides on the harmonic peaks instead of averaging peaks and valleys. The order sets the resolution: about sr / (2 f0) of the sound (e.g. 150 for a voice at 44.1 kHz) follows the formants without rippling at the harmonics}
\fn{(flatten-spectrum mags order)}{the magnitudes divided by their envelope: the fine structure alone}
\fn{(impose-envelope mags source order)}{mags reshaped to carry the spectral envelope of source}
\fn{(stft x n hop)}{=\textgreater{} list of (list re im); (istft frames n hop) =\textgreater{} signal by overlap-add}
\fn{(istft frames n hop)}{the signal back from stft frames by windowed overlap-add}
\fn{(stft-magnitudes frames)}{list of positive-frequency magnitude vectors, one per frame}
\fn{(princarg x)}{wrap a phase to (-pi, pi]}
\fn{(opt opts key default)}{the value of key in a list of (list key value) options, or default}
\fn{(ramp v)}{a parameter given as a number or as (list start end) -\textgreater{} (list start end)}
\fn{(fftshift v)}{rotate a vector by half its length (zero-phase windowing)}
\fn{(pvoc-stretch x factor)}{time-stretch by factor with the pitch kept}
\fn{(pvoc-pitch x ratio)}{pitch-shift by ratio at the same length (formants move too)}
\fn{(pvoc-pitch-formant x ratio order)}{pitch-shift keeping the formants (order about sr/100: 400 at 44.1 kHz)}
\fn{(pvoc-formants x ratio order)}{move the formants by ratio, the pitch kept}
\fn{(pvoc-cross x y mode amount order)}{cross synthesis of x by y: mode 1 multiplicative, 2 spectral flattener (y's envelope on x), 3 morphing}
\fn{(robotize x)}{zero phases: a buzz at the frame rate}
\fn{(whisperize x)}{random phases: the spectral envelope on noise}
\fn{(denoise x threshold)}{magnitudes below threshold x the peak are zeroed (0.01 - 0.1)}
\fn{(gate-spectrum mags threshold)}{magnitudes below threshold times the frame's peak set to zero}
\fn{(spectral-morph a b t n hop)}{between two sounds of the same length: magnitudes interpolated linearly, phases taken from a below t = 0.5 and from b above}
\fn{(spectral-moment amps freqs order centroid)}{weighted moment of the frequencies about a centroid}
\fn{(spectral-centroid amps freqs)}{amplitude-weighted mean frequency}
\fn{(spectral-spread amps freqs)}{standard deviation around the centroid}
\fn{(spectral-skewness amps freqs)}{asymmetry of the spectrum about its centroid}
\fn{(spectral-kurtosis amps freqs)}{peakedness of the spectrum about its centroid}
\fn{(spectral-flux amps previous)}{rectified increase from the previous frame}
\fn{(spectral-irregularity amps)}{sum of absolute differences between neighbouring bins}
\fn{(spectral-decrease amps)}{weighted decrease relative to the first bin}
\fn{(spectral-flatness amps)}{geometric over arithmetic mean: 1 for noise, 0 for a pure tone}
\fn{(spectral-rolloff amps freqs fraction)}{frequency below which the fraction of the energy lies}
\fn{(hfc amps)}{high-frequency content}
\fn{(energy x)}{rms of a frame; (zcr x) zero-crossing rate per sample}
\fn{(zcr x)}{zero-crossing rate per sample}
\fn{(acf-f0 x sr)}{fundamental by autocorrelation; 0 when no clear peak}
\fn{(biquad type sr f0 q gain-db)}{=\textgreater{} (list b a) of an RBJ biquad; type is one of "lowpass" "highpass" "bandpass" "notch" "peak" "lowshelf" "highshelf"}
\fn{(apply-filter x coeffs)}{run x through (list b a) as returned by biquad}
\fn{(lowpass x sr f0 q)  (highpass x sr f0 q)  (bandpass x sr f0 q)  (notch x sr f0 q)}{one biquad, applied}
\fn{(highpass x sr f0 q)}{a highpass biquad, applied}
\fn{(bandpass x sr f0 q)}{a bandpass biquad, applied}
\fn{(notch x sr f0 q)}{a notch biquad, applied}
\fn{(peak-eq x sr f0 q gain-db)  (lowshelf x sr f0 q gain-db)  (highshelf x sr f0 q gain-db)}{the equaliser biquads, applied}
\fn{(lowshelf x sr f0 q gain-db)}{a low shelf, applied}
\fn{(highshelf x sr f0 q gain-db)}{a high shelf, applied}
\fn{(dc-block x)}{remove the DC offset; (dc-block-r x r) with pole radius r (default 0.995)}
\fn{(dc-block-r x r)}{a DC blocker with pole radius r}
\fn{(reson x sr freq tau)}{two-pole resonator at freq Hz with decay time tau seconds; the output lasts tau seconds, so a short excitation rings out}
\fn{(schroeder-reverb x sr rt60)}{four parallel combs into two allpasses; rt60 is the decay time in seconds (how long the tail takes to fall by 60 dB). The output is the input plus rt60 seconds, so the tail is not cut off.}
\fn{(resample-to x sr-in sr-out)}{}
\fn{(envelope-follow x n)}{rms over hops of n samples, one value per hop}
\fn{(envelope-from-values v hop)}{piecewise-linear signal through successive values, hop samples apart}
